\chapter{Results}
\label{ch:resultsmain}

Every number below is labelled with the strength of evidence behind it in the
chapter it is drawn from, and \textbf{no result before the pilot in
\cref{sec:pilot-results} was validated on a physical Kinova arm carrying a
load}: the great majority of this project is simulated against mock hardware,
using the same kinematic description, inverse-kinematics solver and safety
code as the delivered system, plus recorded data from the physical master and
four physical cameras. \Cref{tab:results-summary} is the map from headline
number to the chapter carrying its full derivation.

\begin{table}[htbp]
  \centering
  \small
  \caption[Headline results and their evidence class]{Every headline result
  in this report, its evidence class, and the chapter with the full
  derivation.}
  \label{tab:results-summary}
  \begin{tabular}{@{}p{0.30\linewidth}p{0.46\linewidth}p{0.16\linewidth}@{}}
    \toprule
    Result & Evidence class & Chapter \\
    \midrule
    Six of fourteen master channels alive; incoherent channels inject
      \SI{54.5}{\milli\metre} of unrequested motion (left arm)
      & recorded hardware data, \num{20430} frames & \cref{ch:results} \\
    Vertical and fore-aft command track correctly; lateral measures
      \SI{137.7}{\degree} where \SI{90}{\degree} was intended, and no axis map
      can fix it & recorded hardware data & \cref{ch:results} \\
    The two arms share no reachable point under the as-built mount (0/63
      cells); a costed mount change raises shared cells to 22/63 and home
      clearance from \SI{0.1610}{\metre} to \SI{0.2202}{\metre} & verified in
      simulation & \cref{ch:workspace} \\
    Per-joint step limiter left the hand \SI{51.3}{\milli\metre} off the
      solved path; the jerk-limited replacement is \SI{0.0}{\milli\metre} off
      it and within the manufacturer's joint-speed limit & verified in
      simulation & \cref{ch:motion} \\
    Every joint parks \SI{0.305}{\degree} short of its target on hardware; a
      one-scalar-per-arm correction removes \SIrange{74}{76}{\percent} of it
      & measured on hardware, replayed & \cref{ch:motion} \\
    Positioning error budget: \SI{18.64}{\milli\metre} to
      \SI{12.91}{\milli\metre} RSS, \SI{33.39}{\milli\metre} to
      \SI{19.99}{\milli\metre} worst case, against a \SI{30}{\milli\metre}
      capture gate & measured/derived & \cref{ch:motion} \\
    Fourteen of fourteen injected faults handled; dead-man now keys on the
      source timestamp, not arrival & verified in simulation &
      \cref{ch:results} \\
    Perception on four physical cameras: wrist-camera support plane at
      \SI{1.43}{\milli\metre} RMS, \SI{74}{\percent} inliers; no grasp closed
      on any camera (narrowest object axis \SI{144}{\milli\metre} against an
      \SI{85}{\milli\metre} jaw) & measured on physical hardware &
      \cref{ch:vision} \\
    Task set: \num{13927} inverse-kinematics calls, zero failures, zero
      \SI{150}{\milli\metre} clearance breaches, reproduced field for field on
      re-run & verified in simulation & \cref{ch:results} \\
    \bottomrule
  \end{tabular}
\end{table}

\section{The command mapping and its safety layer}

Channel health, computed over \emph{distinct} sensor updates rather than
transmitted rows (roughly \SIrange{82}{88}{\percent} of consecutive rows
repeat, which had previously been misread as dead channels), found six of
fourteen channels alive. Replaying the recorded incoherent traces with the
operator's input held constant shows those channels alone inject
\SI{54.5}{\milli\metre} of commanded motion on the left arm and
\SI{0.014}{\milli\metre} on the right -- motion nobody asked for, on a master
whose whole reach is \SI{272}{\milli\metre}. Regressing true reach on the
surviving channels shows the left arm retains no observable radial dimension
at all ($R^2=\num{0.133}$) while the right retains it almost completely
($R^2=\num{0.986}$), which names the repair precisely as two channels on the
left arm and predicts that repairing any other channel would not restore
reach.

Vertical and fore-aft command axes track correctly on both arms. The lateral
axis does not: two master motions intended \SI{90}{\degree}
apart arrive \SI{137.7}{\degree} apart at the commanded end effector, and
because a rotation preserves angles, no axis remapping or recalibration can
turn a non-orthogonal measured triad into an orthogonal one. The cause is
mechanical -- the first joint's axis does not tilt as the arm raises, so its
effect on the tip scales with how bent the arm already is -- and is recorded
as a geometric limitation of the sensing choice rather than a calibration
defect.

The safety layer handled fourteen of fourteen injected faults (channel
dropout, board disconnect and reconnect, emergency stop mid-motion,
wrong-object latching, stale detections, session loss, clock
non-monotonicity and more). Two needed several attempts, and both were the
same class of error: a mechanism keyed on the wrong observable. The dead-man
did not trip on ten seconds of total master silence because the master
substitutes stale data and keeps publishing at full rate; keying the
dead-man on the data's own source timestamp fixed it. An earlier version of
the same emergency-stop path also called a blocking real-driver shutdown
first, adding a measured \SI{4}{\second} between the decision to stop and the
arm actually stopping; the fix parks the arm first, from local publishes
only.

\section{Workspace and motion}

Over a 63-cell grid spanning the frontal working volume with both arms
verified at the home pose, \textbf{no cell was reachable by both arms}
simultaneously, and the result survives freeing the wrist yaw and adding
solver retries. A parametric mount sweep shows this is not intrinsic: moving
the mounts inboard, forward and down with additional forward tilt raises
shared cells to 22 of 63 within a buildable \SI{0.20}{\metre} separation.
That mount change was subsequently applied, raising clearance at the home
pose from \SI{0.1610}{\metre} (a reading of the mount itself, since no joint
moves it) to \SI{0.2202}{\metre}, and moving the pick-and-place placement
targets from \SI{0.530}{\metre} off centre to \SI{0.290}{\metre}. Neither
arm crosses the body centre line at any surface height or distance in
\num{3360} swept configurations; the constraint is the wearer's torso, not
their arm posture, since deleting the wearer's arms from the model entirely
moves the boundary by at most \SI{75}{\milli\metre}.

The per-joint step limiter that constituted the whole of teleoperation
motion generation left the hand \SI{51.3}{\milli\metre} (left arm) and
\SI{39.2}{\milli\metre} (right) off the straight line its own
inverse-kinematics solution implied, against a \SI{30}{\milli\metre} grasp
capture gate, and commanded peak joint speeds \num{12.53} times the
manufacturer's declared limit because it never read the limits file. Sweeping
its step size shows the excursion does not tend to zero with a smaller step;
it settles at \SI{68}{\milli\metre}, which is the structural signature of an
algorithm that cannot stay on the line rather than a tuning problem. The
jerk-limited, phase-synchronised replacement is \SI{0.0}{\milli\metre} off
the line and at \num{1.00} times the joint-speed limit.

Replaying recorded joint commands on the physical arm shows every joint
parks \SI{0.305}{\degree} short of its target, on the side it approached
from -- a terminal offset rather than noise, confirmed by propagating the
error through the manipulator Jacobian ($r=\num{0.923}$) rather than by a
proportional-gain model. A single scalar per arm, applied as overshoot in
the direction of travel, removes \SIrange{74}{76}{\percent} of it. Folding
this correction into the positioning error budget, together with a
forward-kinematically derived (rather than asserted) pad offset and a
measured-surface correction to the object-centring bias, takes the combined
budget from \SI{18.64}{} to \SI{12.91}{\milli\metre} in quadrature and from
\SI{33.39}{} to \SI{19.99}{\milli\metre} worst case -- inside the
\SI{30}{\milli\metre} requirement for the first time, \emph{provided} the
still-unmeasured camera-to-wrist transform is small. The first sustained
attempt to drive a physical arm from its own wrist camera measured that
transform directly and found it is not: four poses of a stationary table
disagreed by \SI{8.45}{\degree}, which is \SIrange{34}{59}{\milli\metre} of
position error over the working range (\cref{sec:hardware-session}).

\section{Perception and the task set}

Instrumented at thirty-three stages per camera and pointed at four physical
cameras for the first time (\cref{ch:vision}), the wrist cameras fit a
support plane at \SI{1.43}{\milli\metre} RMS with \SI{74}{\percent} inlier
support at working distance -- within a factor of three of the same code
against a renderer, and the first evidence that the surface-fitting stage
transfers to hardware. No grasp closed on any camera: the narrowest axis of
the nearest object measured \SI{144}{\milli\metre} against an
\SI{85}{\milli\metre} jaw, a statement about the scene rather than the
planner. The same session found, within its first hour, that three files
independently asserted the scene camera's mounting orientation and none had
measured it; the disagreement had silently inverted every earlier scene
figure, and it was invisible to every geometric metric because a consistent
rotation preserves every distance -- it was found only by comparing one
camera's picture against another's.

Within simulation, the full task set solves \num{13927} inverse-kinematics
calls with zero failures and zero breaches of the \SI{150}{\milli\metre}
wearer-clearance floor, and reproduces field for field on re-run. Sweeping
the intent estimator that drives shared autonomy against a mis-declared
table shows why the perception stage matters operationally and not only
geometrically: working from declared coordinates, the estimator becomes
\emph{confidently wrong} \SI{50}{\percent} of the time once the declaration
is one object-pitch out, while working from measured coordinates it stays at
\SI{0.0}{\percent} confidently wrong across the same sweep.

\section{A pilot of direct control against shared autonomy}
\label{sec:pilot-results}

Four operators (P1--P4) completed sixteen recorded sessions under the
protocol of \cref{sec:ethics}: direct VR control and shared-autonomy-assisted
control, each on one of three tasks (object tracking, target reaching,
position matching). Coverage is uneven by design -- not every operator
performed every task under both conditions in the time available -- so the
pooled comparison below mixes tasks together and the paired comparison that
follows it is the part of the data that isolates the condition cleanly.
Full session-by-session data, the coverage table, and all nine figures are
in \cref{app:pilot}.

\begin{table}[htbp]
  \centering
  \small
  \caption[Pilot session summary]{Sixteen pilot sessions, pooled by condition.
  Duration includes any pausing or hesitation and is a session-length rather
  than a pure motor-execution measure. Every e-stop hold in this pilot
  occurred once, near the end of the session, to stop the recording, and is
  excluded as a procedural marker rather than a safety event.}
  \label{tab:pilot-summary}
  \begin{tabular}{@{}lccc@{}}
    \toprule
    Measure ($n=8$ sessions per condition) & Direct & Shared autonomy \\
    \midrule
    Time to finish, mean (s) & 263 & 209 \\
    Hand travel distance, mean (m) & 2.55 & 2.55 \\
    Re-grips per session, mean & 29.4 & 23.9 \\
    Share of session actively steering, mean & \SI{58}{\percent} & \SI{64}{\percent} \\
    Robot's lag behind the operator's hand, mean (mm) & 2.8 & 3.0 \\
    \bottomrule
  \end{tabular}
\end{table}

\begin{figure}[htbp]
  \centering
  \includegraphics[width=0.92\linewidth]{pilot/fig1_time_to_finish.png}
  \caption[Pilot: time to finish, by task and condition]{Time to finish each
  task, pooled across the four operators. Each dot is one recorded session;
  the tick is the mean for that task and condition. Session length includes
  setup and hesitation, so this is the least precise of the pilot's measures;
  \cref{fig:pilot-paired} isolates the condition within a single operator and
  task.}
  \label{fig:pilot-time}
\end{figure}

Pooled across all sixteen sessions (\cref{tab:pilot-summary},
\cref{fig:pilot-time}), shared autonomy sessions averaged shorter completion
times (\SI{209}{\second} against \SI{263}{\second}) and fewer re-grips per
session (\num{23.9} against \num{29.4}), with essentially unchanged hand-travel
distance and a robot lag under \SI{3.1}{\milli\metre} in both conditions.
\Cref{fig:pilot-paired} isolates the six cases in which the \emph{same}
operator performed the \emph{same} task under both conditions: five of six
show a shorter time under shared autonomy, and the median improvement is
consistent with the pooled figure.

\begin{figure}[htbp]
  \centering
  \includegraphics[width=0.72\linewidth]{pilot/fig6_paired_time.png}
  \caption[Pilot: paired time to finish, same operator and task]{The six
  cases where one operator performed the identical task under both
  conditions. This is the fairest comparison available in the pilot: it
  removes both between-operator variance and between-task variance.}
  \label{fig:pilot-paired}
\end{figure}

This pilot measures operator performance only, at pilot scale (four
operators, uneven task coverage), and its e-stop, path-length and lag
measures show no evidence of a safety or responsiveness cost to shared
autonomy over the sessions recorded. It does not include a separate wearer
and therefore cannot speak to the wearer-side hypotheses -- trust and
embodiment under autonomy -- that motivated the two-person design of
\cref{app:study}; \cref{sec:pilot-limits} treats what would be needed to
extend it into that full comparison.
